%============== ESTADO DA ARTE E TRABALHOS RELACIONADOS ==============
\section{Estado da arte e trabalhos relacionados}

Esta seção situa o LibrasAI em relação a trabalhos recentes de reconhecimento
automático de Libras e de outras línguas de sinais que compartilham escolhas
técnicas semelhantes — uso de MediaPipe, arquiteturas CNN/LSTM híbridas e
\textit{datasets} de escala comparável. A Tabela~\ref{tab:relacionados} resume a
comparação.

Fora do contexto brasileiro, bases públicas como a LSA64, da língua de sinais
argentina (RONCHETTI et al., 2016), e a AUTSL, da língua de sinais turca (SINCAN;
KELES, 2020), tornaram-se referências recorrentes para avaliação de modelos de
SLR, o que evidencia a importância de bases próprias e públicas para cada língua
de sinais — lacuna ainda presente no caso da Libras.

\begin{table}[htbp]
\centering
\caption{Trabalhos relacionados e comparação com o LibrasAI.}
\label{tab:relacionados}
\small
\begin{tabular}{@{}p{3.2cm} p{2.6cm} p{6.4cm}@{}}
\toprule
\textbf{Trabalho} & \textbf{Referência} & \textbf{Comparação com o LibrasAI} \\
\midrule
MediaPipe Hands: rastreamento de mãos em tempo real &
Zhang et al., 2020 &
Base do \textit{pipeline} de extração de \textit{landmarks} adotado; não é um
sistema de reconhecimento de língua de sinais. O LibrasAI usa essa solução sobre
um \textit{dataset} próprio de Libras, com normalização espacial pelo pulso. \\
\addlinespace
MINDS-Libras: base de dados e modelos CNN &
Rezende, Almeida e Guimarães, 2021 &
\textit{Dataset} de referência com 1.155 vídeos, 20 sinais e 12 sinalizantes.
Serve de \textit{benchmark} de escala para o protocolo de coleta do LibrasAI (10
participantes, 10 repetições por gesto). \\
\addlinespace
Reconhecimento de sílabas em Libras para HRI &
[PENDENTE: autor], 2026 &
Sistema em tempo real com MediaPipe e CNN, acurácia de 97,4\% — mesmo valor
obtido pelo LibrasAI, ainda que sobre \textit{datasets} e vocabulários distintos,
o que reforça a plausibilidade do resultado preliminar. \\
\addlinespace
Reconhecimento dinâmico via CNN sobre matriz espaço-temporal &
Moreira, 2026 &
Também usa MediaPipe e janela deslizante, mas dispensa redes recorrentes.
Reconhece explicitamente a necessidade de testes com maior diversidade de
usuários para avaliar generalização — limitação também presente no estágio atual
do LibrasAI. \\
\addlinespace
Reconhecimento via \textit{Fast Fourier Transform} e \textit{features}
cinemáticas &
Rego et al., 2025 &
Avaliado no MINDS-Libras, atingiu 92\% de acurácia, superando LSTM (84\%), BiLSTM
(89\%) e CNN (89\%) isoladas — evidência da robustez de abordagens híbridas,
linha seguida pelo LibrasAI. \\
\addlinespace
Reconhecimento de gestos com recorrência e convoluções temporais &
Pigou et al., 2018 &
Demonstra a superioridade de modelos recorrentes (LSTM/GRU) sobre \textit{pooling}
temporal simples para gestos dinâmicos — fundamenta a escolha arquitetural do
LibrasAI. \\
\bottomrule
\end{tabular}
\end{table}
